\chapter{The Vertical Serial Test}

\section{Intention}
The serial test in the previous chapter looked \emph{horizontally} within each
draw: it checks whether \((w_1,w_2,w_3)\) behaves like three independent wheels.

The \textbf{vertical serial test} looks \emph{vertically} through time.  Its
intention is to check, for each wheel separately, whether the wheel's outcome
in one draw is independent of its own outcome in the previous draw.

In plain words:
\begin{quote}
  \emph{Does Wheel 1 ``remember'' what it produced last time?  Does Wheel 2?
  Does Wheel 3?}
\end{quote}

\section{What the Test Checks}
For each wheel, we take the sequence of outcomes across consecutive draws:
\[
  x_1, x_2, x_3, \dots
\]
and count how often we see each \textbf{consecutive pair}
\((x_t, x_{t+1})\).  If draws are independent, then for any two outcomes
\(a,b\),
\[
  P(x_{t}=a, x_{t+1}=b) = P(x_t=a)\,P(x_{t+1}=b).
\]
So the expected pair distribution is the \emph{outer product} of the single-draw
probabilities.

\section{PowerSpin Setup: 25 Categories}
As in the horizontal serial test, we merge the three PowerSpin symbol positions
into one category \texttt{symbol}.  Then each wheel output is one of \(25\)
categories:
\[
  \{1,2,\dots,24,\texttt{symbol}\}.
\]
The per-draw model used by the code is:
\begin{itemize}
  \item each number \(1\)--\(24\) has probability \(1/27\);
  \item \texttt{symbol} has probability \(3/27\).
\end{itemize}

\section{How to Perform the Test (following the code)}
The script you provided implements a very concrete version of the vertical
serial test:

\begin{enumerate}
  \item \textbf{Read each draw's three wheel results.}\\
        For each CSV row where all three wheel results are present, the code
        yields a 3-tuple \((w_1,w_2,w_3)\) for that draw (see
        \texttt{iter\_draw\_wheels}).

  \item \textbf{Normalize categories.}\\
        Each wheel value is normalised to one of the 25 categories:
        integers \(1\)--\(24\), or \texttt{symbol} otherwise.

  \item \textbf{Form consecutive-draw pairs \emph{per wheel}.}\\
        For each wheel independently, the code builds a stream of outcomes
        through time and groups them into \textbf{non-overlapping pairs}:
        \[
          (x_0,x_1),\ (x_2,x_3),\ (x_4,x_5),\ \dots
        \]
        This is important: it does not use overlapping pairs \((x_0,x_1),
        (x_1,x_2), \dots\); it skips one draw after each pair.  If a file ends
        with an unpaired draw, it is paired with the first draw of the next
        file.  (This is implemented by the \texttt{pending} mechanism in
        \texttt{load\_observed\_pairs\_across\_files}.)

  \item \textbf{Count observed pairs.}\\
        For each wheel, we obtain a \(25\times 25\) count matrix
        \(\texttt{observed}[a,b]\), where entry \((a,b)\) records how many times
        the wheel produced outcome \(a\) followed by outcome \(b\) in the next
        draw of the pair.

  \item \textbf{Build the expected pair distribution under independence.}\\
        Create a length-25 probability vector
        \(\texttt{p} = (1/27,\dots,1/27,3/27)\).  Under $H_0$ (independence
        between consecutive draws on the same wheel), the expected probability
        of a pair is
        \[
          \texttt{expected}[a,b] = \texttt{p}[a]\cdot \texttt{p}[b],
        \]
        which is exactly \(\texttt{outer(p,p)}\) (see
        \texttt{expected\_pair\_distribution}).

  \item \textbf{Run a chi-squared test and report one $p$-value per wheel.}\\
        Finally, the code runs a chi-squared goodness-of-fit test
        (Section~\ref{sec:chi-squared-test}) separately for Wheel 1, Wheel 2,
        and Wheel 3.  The output is three $p$-values --- one for each wheel's
        vertical independence check.
\end{enumerate}

\section{Why ``Vertical'' Complements ``Horizontal''}
These two tests look for different failure modes:
\begin{itemize}
  \item the \textbf{horizontal} serial test checks whether the three wheels are
        independent \emph{within the same draw};
  \item the \textbf{vertical} serial test checks whether each wheel is
        independent \emph{across draws in time}.
\end{itemize}

Both can be true or false independently.  A mechanism might have wheels that are
independent within a draw but show time dependence, or vice versa.  That is why
it is valuable to run both.

\section{A Practical Note on Non-Overlapping Pairs}
Using non-overlapping pairs \((x_0,x_1),(x_2,x_3),\dots\) simplifies counting
and reduces dependence between the pairs themselves.

\section{One-Year Result}
Using one year of data (43\,799 non-overlapping pairs per wheel), the vertical
serial test produced the following $p$-values:
\begin{itemize}
  \item Wheel 1: \(p = 0.938107\) \ (pairs \(= 43\,799\))
  \item Wheel 2: \(p = 0.655333\) \ (pairs \(= 43\,799\))
  \item Wheel 3: \(p = 0.741826\) \ (pairs \(= 43\,799\))
\end{itemize}

Wheel 1 is \emph{almost suspicious}: its $p$-value is fairly close to the
\(0.95\) boundary used in the rule-of-thumb interpretation
(Section~\ref{sec:chi-squared-test}).